% !TeX program = XeLaTeX
% !TeX root = ../pujavidhanam.tex
\chapt{श्री-महाप्रदोष-व्रतम्}

\centerline{\small{(मूलम्—श्रीव्रतरत्नाकरः, प्रथमो भागः)}}

\input{purvanga/vighneshwara-puja}

\sect{प्रधान-पूजा — प्रदोष-सांब-सदाशिव-पूजा}

\twolineshloka*
{शुक्लाम्बरधरं विष्णुं शशिवर्णं चतुर्भुजम्}
{प्रसन्नवदनं ध्यायेत् सर्वविघ्नोपशान्तये}

\dnsub{सङ्कल्पः}

ममोपात्त-समस्त-दुरित-क्षयद्वारा श्री-परमेश्वर-प्रीत्यर्थं शुभे शोभने मुहूर्ते अद्य ब्रह्मणः
द्वितीयपरार्धे श्वेतवराहकल्पे वैवस्वतमन्वन्तरे अष्टाविंशतितमे कलियुगे प्रथमे पादे
जम्बूद्वीपे भारतवर्षे भरतखण्डे मेरोः दक्षिणे पार्श्वे शकाब्दे अस्मिन् वर्तमाने व्यावहारिकाणां
प्रभवादीनां षष्ट्याः संवत्सराणां मध्ये \blank\see{app:samvatsara_names} नाम संवत्सरे
\blank\see{app:masa_names}-मासे \blank{}-पक्षे त्रयोदश्यां शुभतिथौ \blank{} वासरयुक्तायां
प्रदोषकाले \blank\see{app:nakshatra_names} नक्षत्र \blank\see{app:yoga_names}-योग
\blank\see{app:karanam_names}-करण-युक्तायां च एवं गुण विशेषण विशिष्टायाम्
अस्यां त्रयोदश्यां शुभतिथौ
अस्माकं सहकुटुम्बानां क्षेमस्थैर्य-धैर्य-वीर्य-विजय-आयुरारोग्य-ऐश्वर्याभिवृद्ध्यर्थम्
अपमृत्युपरिहारार्थं सर्वारिष्टनिवृत्त्यर्थं धर्मार्थकाममोक्ष\-चतुर्विधफलपुरुषार्थसिद्ध्यर्थं
पुत्रपौत्राभि\-वृद्ध्यर्थम् इष्टकाम्यार्थसिद्ध्यर्थं
मम इहजन्मनि पूर्वजन्मनि जन्मान्तरे च सम्पादितानां ज्ञानाज्ञानकृतमहा\-पातकचतुष्टय-व्यतिरिक्तानां
रहस्यकृतानां प्रकाशकृतानां सर्वेषां पापानां सद्य अपनोदनद्वारा
सकल-पापक्षयार्थं श्री-सांब-परमेश्वर-वरप्रसादसिद्ध्यर्थं
यथाशक्ति-ध्यानावाहनादिषोडशोपचारैः
श्री-महाप्रदोष-व्रतांगत्वेन श्री-सांब-सदाशिव-पूजां करिष्ये। तदङ्गं कलशपूजां च करिष्ये।

\input{purvanga/aasana-puja}

\input{purvanga/ghanta-puja}

\begingroup
\renewcommand{\devaName}{देव}
\input{purvanga/kalasha-puja}
\endgroup

\input{purvanga/aatma-puja}

\input{purvanga/pitha-puja}

\input{purvanga/guru-dhyanam}

\sect{षोडशोपचार-पूजा}
\renewcommand{\devAya}{श्री-सांब-सदाशिवाय नमः,}
\begin{center}

\twolineshloka*
{ध्यायेच्चैतन्यमात्मानं सर्वज्ञं सर्वकारणम्}
{व्याघ्रचर्माम्बरधरं नीलग्रीवं त्रिलोचनम्}
\textbf{अस्मिन् बिम्बे/प्रतिमायां/चित्रपटे श्री-सांब-सदाशिवं ध्यायामि।}
\medskip

\twolineshloka*
{आगच्छ देवदेवेश महादेवाभयङ्कर}
{गृहाण सह पार्वत्या त्वं च पूजां मया कृताम्}
\textbf{\devAya{} आवाहयामि।}
\medskip

\twolineshloka*
{विचित्ररत्नरचितं दिव्यास्तरणसंयुतम्}
{स्वर्णसिंहासनं चारु गृहाण सुरपूजित}
\textbf{\devAya{} आसनं समर्पयामि।}
\medskip

\twolineshloka*
{पाद्यं ददामि भगवन् भक्त्या तुभ्यं मया प्रभो}
{त्राहि मां सर्वलोकेश मम पापं च नाशय}
\textbf{\devAya{} पाद्यं समर्पयामि।}
\medskip

\twolineshloka*
{पञ्चवक्राय देवाय पञ्चाक्षरस्वरूपिणे}
{नानापरिमलैर्युक्तं तुभ्यमर्घ्यं ददाम्यहम्}
\textbf{\devAya{} अर्घ्यं समर्पयामि।}
\medskip

\twolineshloka*
{कर्पूरचन्दनैर्युक्तं निर्मलं तु कुशोदकम्}
{चन्द्रमौले मया दत्तं गृहाणाचमनीयकम्}
\textbf{\devAya{} आचमनीयं समर्पयामि।}
\medskip

\twolineshloka*
{शूलिने च नमस्तुभ्यं शङ्कराय नमोऽस्तु ते}
{मध्वाज्यदधिसंयुक्तं मधुपर्कं च गृह्यताम्}
\textbf{\devAya{} मधुपर्कं समर्पयामि।}
\medskip

\twolineshloka*
{मध्वाज्यशर्करायुक्तं फलक्षीरसमन्वितम्}
{पञ्चामृतं प्रदास्यामि स्नानं स्वीकुरु शङ्कर}
\textbf{\devAya{} पञ्चामृतस्नानं समर्पयामि।}
\medskip

\twolineshloka*
{बालेन्दुशेखरेशान सोमसूर्याग्निलोचन}
{गन्धोदकं मयाऽऽनीतं स्नानार्थं प्रतिगृह्यताम्}
\textbf{\devAya{} गन्धोदकस्नानं समर्पयामि।}
\medskip

\twolineshloka*
{नमः कैलासवासाय नमः कालान्तकाय च}
{तोयं शुद्धं मया दत्तं स्नानार्थं प्रतिगृह्यताम्}
\textbf{\devAya{} शुद्धोदकस्नानं समर्पयामि। स्नानानन्तरम् आचमनीयं समर्पयामि।}
\medskip

\twolineshloka*
{दिगम्बराय देवाय नमस्ते कृत्तिवाससे}
{वस्त्रद्वयं दुकूलञ्च समाच्छादय शङ्कर}
\textbf{\devAya{} वस्त्रं समर्पयामि।}
\medskip

\twolineshloka*
{ब्रह्मसूत्रमिदं ब्रह्मन् त्रिगुणं त्रिगुणात्मकम्}
{वाञ्छितार्थप्रसिद्ध्यर्थम् उपवीतं च गृह्यताम्}
\textbf{\devAya{} यज्ञोपवीतं समर्पयामि।}
\medskip

\twolineshloka*
{किरीटं कटकं चैव वलयं हारकुण्डलम्}
{गृहाणाभरणं शम्भो शरणागतवत्सल}
\textbf{\devAya{} आभरणानि समर्पयामि।}
\medskip

\twolineshloka*
{श्रीगन्धं कुङ्कुमोपेतं कर्पूरेण समन्वितम्}
{गृहाण सर्वदेवेश सद्योजात नमो नमः}
\textbf{\devAya{} गन्धान् धारयामि।}
\medskip

\twolineshloka*
{गिरिधन्वन् गिरिपते गिरिजापतये नमः}
{पुत्रपौत्राभिवृद्ध्यर्थम् अक्षतांश्च गृहाण भोः}
\textbf{\devAya{} अक्षतान् समर्पयामि।}
\medskip

\twolineshloka*
{नानाविधानि पुष्पाणि चम्पकादीनि सुव्रत}
{मया दत्तानि सङ्गृह्य पुत्रपौत्रान् प्रवर्धय}
\textbf{\devAya{} पुष्पैः पूजयामि।}

\end{center}

\dnsub{अङ्ग-पूजा}
\begin{longtable}{ll@{~---~}l@{~नमः}l@{~पूजयामि}}
    १. & शिवाय नमः & शिवायै & पादौ\\
    २. & शर्वाय नमः & शर्वाण्यै & गुल्फौ\\
    ३. & ईश्वराय नमः & ईश्वर्यै & जानुनी\\
    ४. & रुद्राय नमः & रुद्राण्यै & ऊरू\\
    ५. & परमात्मने नमः & परमेश्वर्यै & कटिं\\
    ६. & गौरीपतये नमः & गौर्यै & जघनं\\
    ७. & शङ्कराय नमः & शङ्कर्यै & नाभिं\\
    ८. & व्याघ्रचर्मधराय नमः & पीताम्बरधारिण्यै & उदरं\\
    ९. & जगदीश्वराय नमः & जगदीश्वर्यै & हृदयं\\
    १०. & महेश्वराय नमः & महेश्वर्यै & वक्षःस्थलं\\
    ११. & भवाय नमः & भवान्यै & स्तनौ\\
    १२. & महादेवाय नमः & महादेव्यै & स्कन्धौ\\
    १३. & त्र्यम्बकाय नमः & त्र्यम्बकायै & बाहून्\\
    १४. & त्रिपुरारये नमः & त्रिपुरसुन्दर्यै & हस्तान्\\
    १५. & शूलपाणये नमः & शूलपाण्यै & कण्ठं\\
    १६. & कालकण्ठाय नमः & कम्बुकण्ठ्यै & चुबुकं\\
    १७. & भगवते नमः & भगवत्यै & मुखं\\
    १८. & पञ्चवक्राय नमः & सुवक्रायै & नेत्राणि\\
    १९. & त्रिनेत्राय नमः & इन्दीवराक्ष्यै & कर्णौ\\
    २०. & सर्पकुण्डलधराय नमः & रत्नताटङ्कधारिण्यै & नासिकां\\
    २१. & चम्पकनासाय नमः & सुनासायै & ललाटं\\
    २२. & फालनेत्राय नमः & कस्तूरीतिलकायै & शिरः\\
    २३. & जटाधराय नमः & पूर्णकुन्तलधरायै & जटाकलापं\\
    २४. & चन्द्रशेखराय नमः & चन्द्रशेखर्यै & सर्वाण्यङ्गानि\\
\end{longtable}

\needspace{3em}
\begingroup
\setlength{\columnseprule}{1pt}
\let\chapt\sect
\input{../namavali-manjari/100/Shiva_108.tex}

\endgroup

\textbf{\devAya{} नानाविध\-परिमल\-पत्रपुष्पाणि समर्पयामि।}

\sect{उत्तराङ्ग-पूजा}
\begin{center}

\twolineshloka*
{दशाङ्गं च पटीरं च एलागुग्गुलुसंयुतम्}
{धूपं गृहाण देवेश विरूपाक्ष नमोऽस्तु ते}
\textbf{\devAya{} धूपमाघ्रापयामि।}
\medskip

\twolineshloka*
{इन्द्वर्कवह्निनेत्राय पुरत्रयमथेति च}
{घृतवर्तिसुसंयुक्तं दीपोऽयमवलोक्यताम्}
\textbf{\devAya{} अलङ्कारदीपं सन्दर्शयामि।}
\medskip

\twolineshloka*
{शुद्धान्नं पायसादीनि नानाशाकयुतानि च}
{षड्रसादीनि देवेश भुक्त्वा शं कुरु मे सदा}
\textbf{\devAya{} नैवेद्यं निवेदयामि। निवेदनानन्तरम् आचमनीयं समर्पयामि।}
\medskip

\twolineshloka*
{नमस्ते शान्तमनसे सोमनाथाय शम्भवे}
{एलालवङ्गकर्पूरं ताम्बूलं प्रतिगृह्यताम्}
\textbf{\devAya{} ताम्बूलं समर्पयामि।}
\medskip

\twolineshloka*
{चन्द्रादित्यौ च धरणिर्विद्युदग्निस्त्वमेव च}
{त्वमेव सर्वज्योतींषि भज नीराजनं प्रभो}
\textbf{\devAya{} कर्पूरनीराजनं दर्शयामि। कर्पूरनीराजनानन्तरम् आचमनीयं समर्पयामि।}
\medskip

\twolineshloka*
{यो वेदादौ स्वरः प्रोक्तो वेदान्ते च प्रतिष्ठितः}
{तस्य प्रकृतिलीनस्य यः परः स महेश्वरः}
\textbf{\devAya{} मन्त्रपुष्पं समर्पयामि।}
\medskip

\twolineshloka*
{नानारत्नसमायुक्तं वज्रनालसमन्वितम्}
{मुक्ताकेसरसंयुक्तं स्वर्णपुष्पं ददाम्यहम्}
\textbf{\devAya{} सुवर्णपुष्पं समर्पयामि।}
\medskip

\twolineshloka*
{चन्द्रशेखर भूतेश त्रिलोचन वृषध्वज}
{प्रदक्षिणं करिष्यामि प्रकृष्टफलसिद्धये}
\textbf{\devAya{} प्रदक्षिणं समर्पयामि।}
\medskip

\twolineshloka*
{नमस्ते देवदेवेश सृष्टिस्थित्यन्तहेतवे}
{सोमसूर्याग्निनेत्राय नमः सोमार्धमौलये}
\textbf{\devAya{} नमस्कारान् समर्पयामि।}
\medskip

\twolineshloka*
{ऋणपातकदौर्भाग्यदारिद्र्यविनिवृत्तये}
{अशेषाघविनाशाय प्रसीद मम शङ्कर}

\twolineshloka*
{दुःखशोकाभिसन्तप्तं संसारभयपीडितम्}
{बहुपापकृतं दीनं पाहि मां वृषवाहन}
\textbf{\devAya{} प्रार्थनाः समर्पयामि।}

\end{center}

\dnsub{प्रतियाम-अर्घ्यम्}

ममोपात्त-समस्त-दुरित-क्षयद्वारा श्री-सांब-परमेश्वर-प्रीत्यर्थं प्रदोषकाले प्रतियाम-पूजान्ते
क्षीरार्घ्यप्रदानम् उपायनदानं च करिष्ये॥

\twolineshloka*
{कुबेरमित्र देवेश भूतेश त्रिपुरान्तक}
{पार्वतीहृदयानन्द प्रथमार्घ्यं गृहाण भोः}
\textbf{परमेश्वराय नमः इदमर्घ्यम् इदमर्घ्यम् इदमर्घ्यम्।}

\twolineshloka*
{उमाकान्त नमस्तुभ्यं जटामकुटमण्डित}
{चन्द्रमौले त्रिनेत्र त्वं गृहाणार्घ्यं द्वितीयकम्}
\textbf{परमेश्वराय नमः इदमर्घ्यम् इदमर्घ्यम् इदमर्घ्यम्।}

\twolineshloka*
{कैलासवासिन् देवेश भुवनत्रयपालक}
{पञ्चब्रह्मस्वरूप त्वं गृहाणार्घ्यं तृतीयकम्}
\textbf{परमेश्वराय नमः इदमर्घ्यम् इदमर्घ्यम् इदमर्घ्यम्।}

\twolineshloka*
{समुद्रमथनोद्भूत चण्डहालाहलाशन}
{तुरीयार्घ्यं गृहाणेदं मया दत्तं दयानिधे}
\textbf{शिवाय नमः इदमर्घ्यम् इदमर्घ्यम् इदमर्घ्यम्।}

(एतानि चत्वारि अर्घ्याणि प्रतियामम् एकैकं दद्यात्। उत्तराणि सर्वेषु यामेषु।)

\twolineshloka*
{अम्बिकायै नमस्तुभ्यं नमस्ते देवि पार्वति}
{गृहाणार्घ्यं मया दत्तं सर्वसिद्धिप्रदा भव}
\textbf{पार्वत्यै नमः इदमर्घ्यम्।}

\twolineshloka*
{सुब्रह्मण्य महाभाग कार्तिकेय सुरेश्वर}
{इदमर्घ्यं प्रदास्यामि सुप्रीतो भव सर्वदा}
\textbf{सुब्रह्मण्याय नमः इदमर्घ्यम्।}

\twolineshloka*
{नन्दिकेश महाभाग शिवध्यानपरायण}
{इदमर्घ्यं प्रदास्यामि सुप्रीतो भव सर्वदा}
\textbf{नन्दिकेश्वराय नमः इदमर्घ्यम्। चण्डिकेश्वराय नमः इदमर्घ्यम्।}

अनेनार्घ्यप्रदानेन सांबपरमेश्वरः प्रीयताम्॥ अथ उपायनदानं कुर्यात्॥

\hiranyagarbha

श्रीमहाप्रदोषव्रतकाले अस्मिन् मया क्रियमाण-श्रीसांबपरमेश्वरपूजायां
यद्देयमुपायनदानं तत्प्रत्याम्नायार्थं हिरण्यं श्रीसांबपरमेश्वरप्रीतिं
कामयमानः मनसोद्दिष्टाय ब्राह्मणाय तुभ्यमहं सम्प्रददे नमः न मम।
अनया पूजया श्रीसांबपरमेश्वरः प्रीयताम्।

\dnsub{उद्यापनम्}

एवं प्रतियामपूजां कृत्वा अन्ते उपायनानि दत्त्वा रात्रौ जागरणं कृत्वा कथां च श्रुत्वा,
प्रभाते नित्यकर्माणि परिसमाप्य, आचार्यादि-ऋत्विग्वरणं कृत्वा सद्योजातरूपेण आचार्यं,
वामदेव-अघोर-तत्पुरुष-ईशान-रूपेण च क्रमशः अन्यान् ब्राह्मणान् वृत्वा, शुचौ समे देशे
चतुरश्रं स्थण्डिलं कल्पयित्वा तन्मध्ये व्रीहि-तण्डुलोपरि पद्मं लिखित्वा तदुपरि प्रागग्र-दर्भान्
संस्तीर्य तेषु मध्ये एकं परितश्च चतुरः कुम्भान् प्रतिष्ठाप्य वस्त्रादिभिः अलङ्कृत्य
प्रतिमाः स्थापयित्वा, पुण्याहवाचनं कृत्वा सर्वोपकरणानि प्रोक्ष्य, "तत्तत्प्रतिमापूजां करिष्ये,
तदङ्गं कलशपूजां च करिष्ये" इति सङ्कल्प्य कलशपूजां कृत्वा, सर्वे कुम्भं स्पृष्ट्वा
"आपो वा इदꣳ सर्वम्" इत्यादि जपित्वा, कुम्भेषु "इमं मे वरुण" "तत्त्वायामि" इति द्वाभ्यां
वरुणम् आवाह्य, प्रतिमासु ईशानादीन् आवाह्य प्राणप्रतिष्ठां कुर्यात्।

\twolineshloka*
{आवाहयेदुमाकान्तं सद्योजातादिमन्त्रकैः}
{मुक्ताप्रवालसहिते जाम्बूनदपरिष्कृते}

\twolineshloka*
{उपविष्टमुमाकान्तं स्फाटिके वृषवाहने}
{गङ्गाधरं चन्द्रमौलिं व्यालयज्ञोपवीतिनम्}
\textbf{अस्यां प्रतिमायाम् ईशानम् आवाहयामि।}

(एवमेव सर्वासु प्रतिमासु तत्तन्मन्त्रैः तत्तद्देवताम् आवाह्य प्राणप्रतिष्ठां कृत्वा,
आसनादि-उपचारैः पूजयेत्—)

\twolineshloka*
{अनेकरत्नसंयुक्तं मुक्तामणिविभूषितम्}
{ध्यात्वा तेनैव रूपेण आसनं परिकल्पये}
\textbf{ईशानाय नमः आसनं समर्पयामि।}

\twolineshloka*
{पाद्यं ददामि भगवन् भक्त्या तुभ्यं मया प्रभो}
{त्राहि मां सर्वलोकेश मम पापं च नाशय}
\textbf{ईशानाय नमः पाद्यं समर्पयामि।}

\twolineshloka*
{पञ्चवक्राय देवाय पञ्चाक्षरस्वरूपिणे}
{नानापरिमलैर्युक्तं तुभ्यमर्घ्यं ददाम्यहम्}
\textbf{ईशानाय नमः अर्घ्यं समर्पयामि। एवं शुद्धोदकस्नान-वस्त्र-उपवीत-गन्ध-अक्षत-पुष्पादि-सर्वोपचारान् समर्पयामि।}

एवं होमशेषं परिसमाप्य, बिल्व-समित्-तिल-अन्नाज्यैः त्र्यम्बकमन्त्रेण अष्टोत्तरशतं हुत्वा,
पुनः पूजां कृत्वा कुम्भान् उद्वास्य, तत्तीर्थेन "आपो हि ष्ठा" इत्यादि-मन्त्रैः स्नापयित्वा
वस्त्र-पुण्ड्रादीन् धृत्वा तीर्थं प्राश्य, अनन्तरं कलश-वस्त्र-प्रतिमा-दानानि करिष्ये इति सङ्कल्प्य
प्रतिमादानानि दशदानादीनि च कृत्वा, पार्वतीपरमेश्वरं सम्पूज्य ब्राह्मणान् भोजयित्वा
आशिषो वाचयित्वा बन्धुभिस्सह भोजयेत्॥

॥ इति श्रीव्रतरत्नाकरे महाप्रदोषव्रतोद्यापनविधिः सम्पूर्णः ॥

\sect{प्रार्थना}

    \input{../stotra-sangrahah/stotras/shiva/ShivaPanchaksharaStotram.tex}

\sect{अपराध-क्षमापनम्}

\twolineshloka*
{यस्य स्मृत्या च नामोक्त्या तपः पूजाक्रियादिषु}
{न्यूनं सम्पूर्णतां याति सद्यो वन्दे उमापतिम्}

अनया पूजया श्री-सांब-सदाशिवः प्रीयताम्॥

\fourlineindentedshloka*
{कायेन वाचा मनसेन्द्रियैर्वा}
{बुद्‌ध्याऽऽत्मना वा प्रकृतेः स्वभावात्}
{करोमि यद्यत् सकलं परस्मै}
{नारायणायेति समर्पयामि}

\centerline{ॐ तत्सद्ब्रह्मार्पणमस्तु।}

\closesub

\ifbool{katha}{\input{kathas/pradosha-katha}}{}

\closesection
